\chapter{练习路线图}

\section{第一阶段：写小函数但不跳过边界}

第一阶段不要急着写大型项目。目标是把函数输入、输出和错误路径写清楚。建议练习：

\begin{itemize}
  \item 字符串修剪：实现 \code{trim}、\code{split}、\code{starts_with}。
  \item 数字解析：用 \code{std::from_chars} 解析整数，拒绝半截成功。
  \item 成绩统计：空输入、单元素、多元素都要测试。
  \item 用户名校验：长度、字符集合、空字符串都要覆盖。
\end{itemize}

每个练习都要写至少五个测试。不要只测普通输入；普通输入通常最不容易暴露问题。

\section{第二阶段：写小类型并保护不变量}

这一阶段练习类设计：

\begin{itemize}
  \item \code{Money}：用整数分表示金额，禁止把数量和金额混加。
  \item \code{Date}：构造时校验年月日。
  \item \code{TemperatureRange}：保证最低温不高于最高温。
  \item \code{Percent}：范围限制在 \code{0} 到 \code{100}。
\end{itemize}

每个类型都要回答：对象什么时候有效？哪些状态不允许出现？接口是业务动作还是裸 setter？对象是否应该复制？

\section{第三阶段：写资源和容器项目}

这一阶段练习 RAII 和容器选择：

\begin{itemize}
  \item 临时目录：构造创建目录，析构删除目录，禁止复制，允许移动。
  \item LRU 缓存：哈希表加链表，验证映射和链表同步。
  \item 配置加载器：解析、校验、强类型配置分层。
  \item 词频统计器：清洗、统计、排序、输出分开。
\end{itemize}

项目不要一口气写完。先写纯函数，再接 I/O；先测核心逻辑，再测文件入口。

\section{第四阶段：写并发和工程化练习}

并发练习必须小心，不要用睡眠掩盖同步问题：

\begin{itemize}
  \item 线程安全计数器：分别用 \code{mutex} 和 \code{atomic} 实现。
  \item 可关闭队列：支持 \code{push}、\code{pop}、\code{close}。
  \item 任务执行器：使用 \code{jthread} 管理生命周期。
  \item 配置热加载：先解析到临时对象，再原子替换共享配置。
\end{itemize}

每个并发练习都要写出共享状态、不变量、锁保护范围和退出协议。写不出来这些，就先不要写代码。

\section{第五阶段：复盘自己的代码}

每完成一个练习，按下面问题复盘：

\begin{enumerate}
  \item 哪个函数最难测试？为什么？
  \item 哪个类型的不变量最重要？测试有没有覆盖？
  \item 有没有无意复制大对象？
  \item 有没有保存非拥有视图到过长生命周期？
  \item 错误消息能不能帮助用户定位问题？
  \item 如果输入扩大一百倍，瓶颈可能在哪里？
\end{enumerate}

复盘比刷更多语法更重要。高级 \cpp 能力不是知道更多特性，而是能在复杂约束下写出可解释、可测试、可维护的代码。
